% Resumen del mundo «Raíces de Polinomios» (Álgebra, mundo 8: polinomios de grado superior).
% Lecciones alg-23b, alg-92, alg-93. Ejemplos de los videos (la cordillera de Ana, la apuesta de Beto, la de Caro);
% el taller trae ejercicios nuevos (no repite los de la app, del jefe ni de la evaluación).
\documentclass[11pt,a4paper]{article}
\usepackage{zuka-resumen}
\zukacurso{algebra}
\zukamundo{8}{Raíces de Polinomios}{Polinomios de grado superior}

\newcommand{\gr}{\operatorname{gr}}
% fila de una división larga sin la barra vertical
\newcommand{\nb}{\multicolumn{1}{r@{\;}}{}}

\begin{document}

\portada{%
Un polinomio de grado tres dibuja una \clave{cordillera}: sube, baja y cruza la vía del tren (el eje $x$) en sus
\clave{ceros}. En este mundo aprendiste a dividir polinomios, a calcular residuos sin dividir y a encontrar todos
los ceros de un polinomio, incluso los que no se ven en el dibujo. Este resumen reúne las definiciones, los teoremas
y los ejemplos de los videos, y termina con un taller de práctica.}{%
\item \textbf{Nivel 8-1.} División e igualdad de polinomios: dividendo, divisor, cociente y residuo; comparar coeficientes.
\item \textbf{Nivel 8-2.} Teorema del residuo, teorema del factor y división sintética.
\item \textbf{Nivel 8-3.} Ceros de un polinomio: teorema fundamental, multiplicidad, ceros racionales, regla de Descartes y ceros complejos.}

\begin{clavebox}[LAS TRES IDEAS QUE SOSTIENEN TODO EL MUNDO]
\renewcommand{\arraystretch}{1.5}\begin{tabular}{@{}c@{\qquad}c@{\qquad}c@{}}
$P=D\cdot Q+R$ & $R=P(c)$ al dividir entre $x-c$ & $P(c)=0\iff (x-c)\mid P$\\[2pt]
\small toda división & \small el residuo es una altura & \small cada cero aporta un factor
\end{tabular}
\end{clavebox}

%=====================================================================
\nivel{8-1}{División e igualdad de polinomios}

Ana dibuja el perfil de una cordillera con el polinomio
\[ P(x)=x^3-2x^2-5x+6 . \]
La vía del tren va a altura cero y uno de sus cruces es $x=1$. Ana quiere «quitar» ese cruce del polinomio:
¿qué queda si divide entre $x-1$? Se divide igual que con números.

\begin{minipage}[c]{0.5\linewidth}
\textbf{Con números.} $47$ caramelos en bolsas de $5$ dan $9$ bolsas y sobran $2$:
\[ 47 = 5\cdot 9 + 2,\qquad 0\le 2<5 . \]
Lo que sobra siempre es \emph{menor} que una bolsa; si fuera mayor, cabría otra bolsa.
\end{minipage}\hfill
\begin{minipage}[c]{0.46\linewidth}\centering
\begin{tikzpicture}[x=0.5cm,y=0.5cm]
  \foreach \b in {0,...,8} {
    \draw[zmute,rounded corners=1.5pt,fill=zsky!12] (\b*1.15-0.1,-0.1) rectangle (\b*1.15+0.9,4.9);
    \foreach \k in {0,...,4} \fill[zsky] (\b*1.15+0.4,\k*0.95+0.38) circle (0.27);}
  \fill[zcoral] (11.4,0.38) circle (0.27); \fill[zcoral] (11.4,1.33) circle (0.27);
  \node[font=\footnotesize,text=zmute] at (4.6,-0.85) {$9$ bolsas de $5$};
  \node[font=\footnotesize,text=zcoral] at (11.4,-0.85) {sobran $2$};
\end{tikzpicture}
\end{minipage}

\begin{definicion}[Algoritmo de la división]
Si $P(x)$ y $D(x)$ son polinomios y $D\neq 0$, existen polinomios únicos $Q(x)$ (\clave{cociente}) y $R(x)$
(\clave{residuo}) tales que
\[ P(x)=D(x)\cdot Q(x)+R(x),\qquad R=0\ \text{ o }\ \gr R<\gr D . \]
$P$ es el \clave{dividendo} y $D$ el \clave{divisor}. Si $R=0$, la división es \clave{exacta} y $D$ es un
\clave{factor} de $P$. Además, $\gr Q=\gr P-\gr D$.
\end{definicion}

\begin{ejemplo}[La cordillera de Ana entre $x-1$]
En cada paso se divide el término principal que queda entre $x$, se multiplica por el divisor y se resta:
\[
\renewcommand{\arraystretch}{1.15}
\begin{array}{r@{\;}|rrrr}
\nb & & x^2 & -x & -6\\ \cline{2-5}
x-1 & x^3 & -2x^2 & -5x & +6\\
\nb & -x^3 & +x^2 & & \\ \cline{2-3}
\nb & & -x^2 & -5x & \\
\nb & & +x^2 & -x & \\ \cline{3-4}
\nb & & & -6x & +6\\
\nb & & & +6x & -6\\ \cline{4-5}
\nb & & & & 0
\end{array}
\qquad
\begin{aligned}
&\tfrac{x^3}{x}=x^2\\[2pt]
&\tfrac{-x^2}{x}=-x\\[2pt]
&\tfrac{-6x}{x}=-6\\[2pt]
&\text{residuo } 0
\end{aligned}
\]
El residuo es $0$: $x-1$ cabe justo, es un factor, y
\[ x^3-2x^2-5x+6=(x-1)\left(x^2-x-6\right)=(x-1)(x-3)(x+2). \]
La parábola que queda corta la vía en $-2$ y en $3$: son los otros dos cruces de la cordillera.
\end{ejemplo}

\begin{ojo}[Ordena y completa]
Antes de dividir, ordena los dos polinomios de mayor a menor grado y escribe con coeficiente $0$ los términos que
faltan: $x^3-7x+6$ se divide como $x^3+0x^2-7x+6$. Si no dejas el hueco, las columnas se corren.
\end{ojo}

\subsection{Igualdad de polinomios}
\begin{propiedad}[Dos polinomios son iguales si y solo si tienen los mismos coeficientes]
\[ a_nx^n+\dots+a_1x+a_0=b_nx^n+\dots+b_1x+b_0\ \text{ para todo }x \iff a_k=b_k\ \text{ para cada }k . \]
\end{propiedad}
Esto da un camino para dividir \emph{sin dividir}: se escribe la forma del resultado con letras y se comparan coeficientes.
\begin{ejemplo}[El mismo cociente, comparando coeficientes]
Como $\gr Q=3-1=2$ y el residuo es una constante, se busca $P(x)=(x-1)\left(x^2+ax+b\right)+r$. Al multiplicar:
\[ (x-1)\left(x^2+ax+b\right)+r=x^3+(a-1)x^2+(b-a)x+(r-b). \]
Comparando con $x^3-2x^2-5x+6$:
\[ a-1=-2\Rightarrow a=-1,\qquad b-a=-5\Rightarrow b=-6,\qquad r-b=6\Rightarrow r=0 . \]
Mismo cociente $x^2-x-6$ y mismo residuo $0$.
\end{ejemplo}

%=====================================================================
\nivel{8-2}{Residuo, factor y división sintética}

Beto ve la cordillera de Ana entre la niebla y apuesta: «la vía también cruza en $x=2$». Para saberlo habría que
dividir entre $x-2$\dots\ pero el residuo guarda un secreto.

\begin{propiedad}[Teorema del residuo]
El residuo de dividir $P(x)$ entre $x-c$ es $P(c)$.
\end{propiedad}
\textbf{Por qué.} Toda división se escribe $P(x)=(x-c)\,Q(x)+R$, con $R$ constante. Al poner $x=c$ el primer
pedazo vale $0\cdot Q(c)=0$, y queda $P(c)=R$. El residuo es la \clave{altura} del terreno en $x=c$.

\begin{minipage}[c]{0.47\linewidth}
\begin{ejemplo}[La apuesta de Beto]
\begin{align*}
P(2)&=8-8-10+6=-4,\\
P(3)&=27-18-15+6=0 .
\end{align*}
En $x=2$ el terreno está $4$ unidades por debajo de la vía: Beto pierde. En $x=3$ la altura es $0$: ahí sí cruza.
\end{ejemplo}
\end{minipage}\hfill
\begin{minipage}[c]{0.5\linewidth}\centering
\begin{tikzpicture}
\begin{axis}[zplot,width=7.4cm,height=5.6cm,xmin=-3,xmax=4,ymin=-9,ymax=10,
  xtick={-1,2},ytick={-8,-4,4,8},xlabel={$x$},ylabel={$y$},
  xlabel style={at={(ticklabel* cs:1)},anchor=north west},ylabel style={at={(ticklabel* cs:1)},anchor=south west}]
  \addplot[zaccent,domain=-2.55:3.65,samples=120] {x^3-2*x^2-5*x+6};
  \addplot[zcoral,dashed,line width=0.9pt] coordinates {(2,0) (2,-4)};
  \addplot[only marks,mark=*,mark size=2.2pt,zink] coordinates {(-2,0) (1,0) (3,0)};
  \addplot[only marks,mark=*,mark size=2.2pt,zcoral] coordinates {(2,-4)};
  \node[font=\footnotesize,text=zcoral,anchor=west] at (axis cs:2.1,-5.4) {$P(2)=-4$};
  \node[font=\footnotesize,text=zaccent,anchor=west] at (axis cs:0.35,8.2) {$y=P(x)$};
\end{axis}
\end{tikzpicture}

{\small\color{zmute}Los cruces $-2$, $1$ y $3$ y la altura en $x=2$.}
\end{minipage}

\begin{propiedad}[Teorema del factor]
\[ P(c)=0 \iff (x-c)\ \text{es un factor de } P(x). \]
Cada cruce $c$ con la vía da un factor $x-c$, y cada factor $x-c$ da un cruce en $c$.
\end{propiedad}

\subsection{División sintética}
Para dividir entre $x-c$ hay una tabla que divide \emph{sin escribir las $x$}. Se ponen los coeficientes en fila
(con ceros donde falte un término) y $c$ a la izquierda. Se baja el primero; luego, en cada columna, se
\emph{multiplica} lo último de abajo por $c$, se escribe arriba y se \emph{suma}.

\begin{ejemplo}[$P(x)=x^3-2x^2-5x+6$ entre $x-3$]
\[
\renewcommand{\arraystretch}{1.3}
\begin{array}{c|rrrr}
3 & 1 & -2 & -5 & 6\\
  &   & 3 & 3 & -6\\ \hline
  & 1 & 1 & -2 & \fbox{$0$}
\end{array}
\qquad\Longrightarrow\qquad
Q(x)=x^2+x-2,\quad R=0 .
\]
La última casilla es el residuo y las demás son los coeficientes del cociente, que tiene un grado menos. Como
$x^2+x-2=(x+2)(x-1)$, otra vez $P(x)=(x-3)(x+2)(x-1)$.
\end{ejemplo}

\begin{ojo}[El signo de $c$]
Para dividir entre $x+2$ se escribe $c=-2$, porque $x+2=x-(-2)$. Y la división sintética solo sirve con divisores
de la forma $x-c$; para $2x-1$ se usa $c=\tfrac12$ y luego se divide el cociente entre $2$.
\end{ojo}

%=====================================================================
\nivel{8-3}{Ceros de un polinomio}

La cordillera de Caro también es de grado tres, pero cruza la vía una sola vez, en $x=1$:
\[ Q(x)=x^3-3x^2+4x-2 . \]
Y un teorema promete tres ceros. ¿Dónde están los otros dos?

\begin{propiedad}[Teorema fundamental del álgebra]
Todo polinomio de grado $n\ge 1$ tiene exactamente $n$ ceros complejos, si cada uno se cuenta tantas veces como
su multiplicidad. Es decir, se factoriza como
\[ P(x)=a_n\,(x-r_1)(x-r_2)\cdots(x-r_n),\qquad r_k\in\C . \]
\end{propiedad}

\subsection{Multiplicidad: cruzar o rebotar}
\begin{minipage}[c]{0.5\linewidth}
\begin{definicion}[Multiplicidad]
Si $(x-r)^m$ es factor de $P$ pero $(x-r)^{m+1}$ no lo es, $r$ es un cero de \clave{multiplicidad} $m$.
\end{definicion}
En $y=(x+1)(x-2)^2$ el cero $-1$ es simple y el $2$ es \emph{doble}: cuenta dos veces, y $1+2=3$ es el grado.
\begin{itemize}
\item Multiplicidad \textbf{impar}: la gráfica \clave{cruza} el eje.
\item Multiplicidad \textbf{par}: lo \clave{toca y rebota}.
\end{itemize}
\end{minipage}\hfill
\begin{minipage}[c]{0.46\linewidth}\centering
\begin{tikzpicture}
\begin{axis}[zplot,clip=false,width=7cm,height=5.4cm,xmin=-2.3,xmax=3.6,ymin=-6,ymax=7,
  xtick={1,3},ytick={-4,4},xlabel={$x$},ylabel={$y$},
  xlabel style={at={(ticklabel* cs:1)},anchor=north west},ylabel style={at={(ticklabel* cs:1)},anchor=south west}]
  \addplot[zaccent,domain=-1.45:3.25,samples=120] {(x+1)*(x-2)^2};
  \addplot[only marks,mark=*,mark size=2.2pt,zink] coordinates {(-1,0) (2,0)};
  \node[font=\footnotesize,text=zink,anchor=south] at (axis cs:-1.5,0.9) {cruza en $-1$};
  \node[font=\footnotesize,text=zink,anchor=north] at (axis cs:2,-0.9) {rebota en $2$};
\end{axis}
\end{tikzpicture}

{\small\color{zmute}$y=(x+1)(x-2)^2$: cero simple y cero doble.}
\end{minipage}

\subsection{Cómo buscar los ceros}
\begin{propiedad}[Ceros racionales]
Si $P(x)=a_nx^n+\dots+a_1x+a_0$ tiene coeficientes enteros y $\tfrac{p}{q}$ (fracción irreducible) es un cero de
$P$, entonces $p$ divide a $a_0$ y $q$ divide a $a_n$.
\end{propiedad}
\begin{propiedad}[Regla de los signos de Descartes]
El número de ceros reales \emph{positivos} de $P$ es igual al número de cambios de signo entre coeficientes
consecutivos (sin contar los ceros), o es menor que ese número en una cantidad par. Para los ceros
\emph{negativos} se cuentan los cambios de signo de $P(-x)$.
\end{propiedad}
\begin{propiedad}[Ceros conjugados]
Si $P$ tiene coeficientes reales y $a+bi$ es un cero, también lo es $a-bi$: los ceros complejos llegan en pareja.
\end{propiedad}

\begin{ejemplo}[Los ceros perdidos de Caro]
\textbf{Candidatos.} Con $a_0=-2$ y $a_n=1$, los ceros racionales posibles son $\pm1,\ \pm2$.\par
\textbf{Descartes.} Los signos de $Q$ son $+,-,+,-$: tres cambios, así que hay $3$ o $1$ ceros positivos.
En $Q(-x)=-x^3-3x^2-4x-2$ no hay cambios de signo: ningún cero negativo.\par
\textbf{Probar.} $Q(1)=1-3+4-2=0$, pero $Q(2)=8-12+8-2=2\neq0$.\par\smallskip
\begin{minipage}[c]{0.6\linewidth}
\textbf{Dividir entre $x-1$.}
\[
\renewcommand{\arraystretch}{1.3}
\begin{array}{c|rrrr}
1 & 1 & -3 & 4 & -2\\
  &   & 1 & -2 & 2\\ \hline
  & 1 & -2 & 2 & \fbox{$0$}
\end{array}
\]
La parábola $x^2-2x+2$ tiene discriminante $(-2)^2-4\cdot2=-4<0$, así que
\[ x=\frac{2\pm\sqrt{-4}}{2}=1\pm i . \]
Los tres ceros son $1$, $1+i$ y $1-i$:
\[ Q(x)=(x-1)\left(x^2-2x+2\right). \]
\end{minipage}\hfill
\begin{minipage}[c]{0.36\linewidth}\centering
\begin{tikzpicture}[scale=1.1]
  \draw[zline,line width=0.4pt] (-0.8,-1.6) grid (2.3,1.6);
  \draw[->,zink,line width=0.8pt] (-0.9,0) -- (2.5,0) node[right,font=\small] {Re};
  \draw[->,zink,line width=0.8pt] (0,-1.7) -- (0,1.8) node[above,font=\small] {Im};
  \node[font=\footnotesize,text=zmute,below right] at (1,0) {$1$};
  \node[font=\footnotesize,text=zmute,below] at (2,0) {$2$};
  \node[font=\footnotesize,text=zmute,left] at (0,1) {$i$};
  \node[font=\footnotesize,text=zmute,left] at (0,-1) {$-i$};
  \draw[zcoral,dashed,line width=0.8pt] (1,1) -- (1,-1);
  \fill[zink] (1,0) circle (2.4pt);
  \fill[zcoral] (1,1) circle (2.4pt) node[right,font=\small,text=zink] {$1+i$};
  \fill[zcoral] (1,-1) circle (2.4pt) node[right,font=\small,text=zink] {$1-i$};
\end{tikzpicture}

{\small\color{zmute}Los ceros complejos son espejo uno del otro respecto al eje real.}
\end{minipage}
\end{ejemplo}

\begin{ojo}[Lo que dicen y lo que no dicen los teoremas]
Los candidatos racionales son \emph{posibles} ceros: hay que probarlos. Descartes da un máximo que baja de dos en
dos, no un número seguro. Y que la gráfica cruce una sola vez no quiere decir que haya un solo cero: los demás
pueden ser complejos o estar repetidos.
\end{ojo}

%=====================================================================
\clearpage
\repaso
\begin{center}
\begin{tabularx}{\linewidth}{@{}l L@{}}
\toprule
Resultado & Qué dice\\ \midrule
Algoritmo de la división & $P=D\cdot Q+R$, con $R=0$ o $\gr R<\gr D$; además $\gr Q=\gr P-\gr D$\\
Igualdad de polinomios & iguales para todo $x$ $\iff$ mismos coeficientes, uno a uno\\
Teorema del residuo & el residuo de $P\div(x-c)$ es $P(c)$\\
Teorema del factor & $P(c)=0\iff(x-c)$ es factor de $P$\\
División sintética & baja, multiplica por $c$, suma; la última casilla es el residuo\\
Teorema fundamental & grado $n\ge1$ $\Rightarrow$ $n$ ceros complejos contados con multiplicidad\\
Multiplicidad & exponente del factor; impar cruza, par rebota\\
Ceros racionales & $\tfrac{p}{q}$ con $p\mid a_0$ y $q\mid a_n$\\
Descartes & ceros positivos: cambios de signo de $P(x)$, o menos en un par; negativos: los de $P(-x)$\\
Conjugados & coeficientes reales: si $a+bi$ es cero, también $a-bi$\\
\bottomrule
\end{tabularx}
\end{center}

\begin{ojo}[Errores frecuentes]
\begin{itemize}
\item No escribir los coeficientes $0$ de los términos que faltan, en la división larga o en la sintética.
\item Usar $c=2$ para dividir entre $x+2$ (es $c=-2$).
\item Restar mal en la división larga: se cambia el signo de \emph{toda} la fila antes de sumar.
\item Dar un residuo de grado mayor o igual que el divisor: todavía se puede seguir dividiendo.
\item Tomar los candidatos racionales como ceros seguros, o leer Descartes como un número exacto.
\item Olvidar los ceros complejos o la multiplicidad al contar los $n$ ceros.
\end{itemize}
\end{ojo}

%=====================================================================
\taller
\begin{ejercicios}
\item Divide $\left(2x^3+3x^2-4x+5\right)\div(x+2)$ con división larga y escribe el resultado en la forma
  $P=D\cdot Q+R$.
\item Halla el cociente y el residuo de $\left(x^4-3x^2+2x-1\right)\div\left(x^2+x-1\right)$. Comprueba que
  $\gr R<\gr D$.
\item Halla $a$, $b$ y $c$ tales que $3x^2-5x+1=a(x-1)^2+b(x-1)+c$ para todo $x$.
\item Sin hacer la división, halla el residuo:
  \begin{ejercicios}
  \item de $x^5-3x^3+x-4$ entre $x+1$;
  \item de $4x^3-6x^2+3$ entre $2x-1$. (Sugerencia: $2x-1$ se anula en $x=\tfrac12$.)
  \end{ejercicios}
\item Determina $k$ para que $x+2$ sea factor de $x^3+kx^2-x-2$ y, con ese valor, factoriza el polinomio por
  completo.
\item Usa división sintética para dividir $x^4-5x^3+5x^2+5x-6$ entre $x-3$. Luego factoriza el polinomio en
  factores lineales.
\item Sea $P(x)=2x^3-x^2-13x-6$.
  \begin{ejercicios}
  \item Usa la regla de Descartes para decir cuántos ceros positivos y cuántos negativos puede tener.
  \item Lista los candidatos a ceros racionales y halla todos los ceros. ¿Concuerda con lo que dijo Descartes?
  \end{ejercicios}
\item Un polinomio $P$ de grado $5$ con coeficientes reales tiene los ceros $2$ (doble), $-1$ y $3i$.
  \begin{ejercicios}
  \item ¿Cuál es el quinto cero? Justifica.
  \item Escribe $P$ factorizado si su coeficiente principal es $1$, y di en qué ceros reales la gráfica cruza el eje y en cuáles rebota.
  \end{ejercicios}
\item La figura muestra la gráfica de un polinomio $f$ de grado $4$ que pasa por el punto $(0,6)$.
  \begin{center}
  \begin{tikzpicture}
  \begin{axis}[zplot,width=8cm,height=5.6cm,xmin=-3.2,xmax=3.8,ymin=-10,ymax=13,
    xtick={-3,-2,-1,1,2,3},ytick={-8,-4,4,8,12},xlabel={$x$},ylabel={$y$},
    xlabel style={at={(ticklabel* cs:1)},anchor=north west},ylabel style={at={(ticklabel* cs:1)},anchor=south west}]
    \addplot[zaccent,domain=-2.85:3.38,samples=140] {0.5*(x+2)^2*(x-1)*(x-3)};
    \addplot[only marks,mark=*,mark size=2pt,zink] coordinates {(0,6)};
  \end{axis}
  \end{tikzpicture}
  \end{center}
  \begin{ejercicios}
  \item Lee los ceros de $f$ y la multiplicidad de cada uno (piensa en cuáles cruza y en cuáles rebota).
  \item Escribe una fórmula para $f$ y usa el punto $(0,6)$ para hallar su coeficiente principal.
  \end{ejercicios}
\item Con una lámina de cartón de $\SI{20}{cm}\times\SI{12}{cm}$ se arma una caja sin tapa: en cada esquina se
  corta un cuadrado de lado $x$ y se doblan los bordes.
  \begin{ejercicios}
  \item Muestra que el volumen es $V(x)=x(20-2x)(12-2x)$ y di para qué valores de $x$ tiene sentido.
  \item ¿Para qué valores de $x$ el volumen es de $\SI{256}{cm^3}$? (Sugerencia: prueba $x=2$ y divide.)
  \end{ejercicios}
\end{ejercicios}

\clearpage
\respuestas
\begin{respuestascortas}
\begin{soluciones}
\item $2x^3+3x^2-4x+5=(x+2)\left(2x^2-x-2\right)+9$.
\item $Q=x^2-x-1$, $R=2x-2$; $\gr R=1<2$.
\item $a=3$, $b=1$, $c=-1$ (el coeficiente de $x$ da $-2a+b=-5$).
\item (a) $P(-1)=-1+3-1-4=-3$. (b) $P\!\left(\tfrac12\right)=\tfrac12-\tfrac32+3=2$.
\item $P(-2)=4k-8=0\Rightarrow k=2$; $x^3+2x^2-x-2=(x+2)(x-1)(x+1)$.
\item Fila: $1,\,-2,\,-1,\,2,\,\fbox{$0$}$, así que $Q=x^3-2x^2-x+2=(x-2)(x-1)(x+1)$; el polinomio es $(x-3)(x-2)(x-1)(x+1)$.
\item (a) $P$: $+,-,-,-$, un cambio: exactamente $1$ positivo. $P(-x)$: $-,-,+,-$, dos cambios: $2$ o $0$ negativos.
  (b) Candidatos $\pm1,\pm2,\pm3,\pm6,\pm\tfrac12,\pm\tfrac32$. $P(3)=0$ y $P=(x-3)(2x+1)(x+2)$: ceros $3,\,-\tfrac12,\,-2$. Sí concuerda.
\item (a) $-3i$, porque con coeficientes reales los complejos vienen en pareja conjugada. (b) $P(x)=(x-2)^2(x+1)\left(x^2+9\right)$; cruza en $-1$ y rebota en $2$.
\item (a) $-2$ doble (rebota), $1$ y $3$ simples (cruza). (b) $f(x)=a(x+2)^2(x-1)(x-3)$ con $f(0)=12a=6$: $a=\tfrac12$.
\item (a) $0<x<6$. (b) $4x(10-x)(6-x)=256\iff x^3-16x^2+60x-64=0$; con $x=2$ queda $x^2-14x+32=0$, $x=7\pm\sqrt{17}$. Sirven $x=\SI{2}{cm}$ y $x=7-\sqrt{17}\approx\SI{2.88}{cm}$ ($7+\sqrt{17}>6$ no sirve).
\end{soluciones}
\end{respuestascortas}

\end{document}
